%%%%%%%%%%%%%%%%%%%%
%
% $Autor: Wings $
% $Datum: 2020-02-24 14:29:03Z $
% $Version: 1791 $
%
%
%%%%%%%%%%%%%%%%%%%





\section{Setting up the Camera}

After starting the system, the Raspberry Pi icon menu can be activated at the top left. There you can go to the settings \menu[>]{Preferences > Raspberry Pi Configuration > Interfaces}. To use the camera, it must then be activated (Enabled). After activating the camera, the computer must be restarted. A dialogue with a camera command line can be opened with \keys{\ctrl + \Alt + T}. There you can make a test shot of an image by entering the following information:

\begin{lstlisting}
raspistill -v -o test.jpg
xdg-open /home/pi/test.jpg
\end{lstlisting}

The picture is then stored in the file \FILE{test.jpg}.


%\begin{notes}
%\item Das Pi-Kameraprogramm des Kameramoduls hilft bei der Durchführung der Echtzeit-Live-Identifikation von Objekten.
%Damit das Pi-Kamera-Modul speziell für die Live-Objektidentifizierung im VNC-Viewer funktioniert, aktivieren Sie den Direktaufnahmemodus des VNC-Viewers.
%\end{notes}
